% 基于上一轮模型设计重建为真实矢量：同一句文本的三个任务。
\begingroup%
% 第1—3章局部矢量图元：单色黄、双色蓝红、三色蓝红黄。
\providecommand{\BasicsFiveSetup}{%
 \colorlet{B5Blue}{FigureBlue}\colorlet{B5Coral}{FigureRed}%
 \colorlet{B5Yellow}{FigureYellow}\definecolor{B5Gray}{HTML}{4C4D4F}%
 \colorlet{B5Ice}{FigureBlueFill}\colorlet{B5Rose}{FigureRedFill}%
 \colorlet{B5Cream}{FigureYellowFill}%
 \tikzset{b5 picture/.style={x=1mm,y=1mm,font=\figurefont\mdseries\fontsize{8.5}{10.5}\selectfont,text=B5Gray,line width=1pt},%
 b5 node/.style={draw=B5Gray,fill=white,line width=1pt,rounded corners=1.5mm,inner sep=3pt,font=\figurefont\mdseries\fontsize{8.5}{10.5}\selectfont,text=B5Gray,align=center},%
 b5 flow/.style={draw=B5Gray,line width=1.1pt,-{Stealth[length=2.2mm,width=1.4mm]}},%
 b5 link/.style={draw=B5Gray,line width=.7pt},%
 b5 feedback/.style={b5 flow,draw=B5Coral,dash pattern=on 2.5mm off 1.5mm}}}%
% 3—3—1网络仅为图标，不指定真实架构；各层用实心蓝、红、黄，连线中性灰。
\providecommand{\BasicsFiveNet}[3]{\begin{scope}[shift={(#1,#2)},scale=#3]%
 \foreach \a in {-3,0,3}{\foreach \b in {-3,0,3}{\draw[b5 link] (0,\a)--(5,\b);}}%
 \foreach \a in {-3,0,3}{\draw[b5 link] (5,\a)--(10,0);}%
 \foreach \a in {-3,0,3}{%
  \draw[draw=B5Gray,line width=.8pt,fill=B5Blue] (0,\a) circle (1.1);%
  \draw[draw=B5Gray,line width=.8pt,fill=B5Coral] (5,\a) circle (1.1);}%
 \draw[draw=B5Gray,line width=.8pt,fill=B5Yellow] (10,0) circle (1.1);%
\end{scope}}%
% 鸟与猫保留教学对象，不生成或模拟真实数据样本。
\providecommand{\BasicsFiveBird}[4]{\begin{scope}[shift={(#1,#2)},scale=#3]%
 \draw[draw=B5Gray,fill=#4,line width=.8pt,line join=round]%
  (-6,-1)--(-10,-2)--(-6,2).. controls (-3,3) and (-2,6) .. (1,6)%
  .. controls (4,6) and (5,4) .. (4,2).. controls (4,-1) and (-1,-3) .. (-6,-1);%
 \draw[draw=B5Gray,fill=white,line width=.7pt] (4,3)--(6,3.8)--(4,4.1)--cycle;%
 \fill[B5Gray] (2.7,4.3) circle (.35);%
 \draw[b5 link] (-5,1).. controls (-1,4) and (1,1) .. (-3,-.5);%
 \draw[b5 link] (-2,-2)--(-2,-4)--(-.7,-4);\draw[b5 link] (0,-1.5)--(0,-4)--(1.3,-4);%
\end{scope}}%
\providecommand{\BasicsFiveCat}[3]{\begin{scope}[shift={(#1,#2)},scale=#3]%
 \draw[draw=B5Gray,fill=white,line width=.8pt] (0,-1) ellipse (3.3 and 4.3);%
 \draw[draw=B5Gray,fill=white,line width=.8pt,line join=round]%
  (-3,3)--(-3,7)--(-1.2,5.8).. controls (0,6.5) and (1,6.5) .. (2,5.8)%
  --(3.3,7)--(3.3,3).. controls (2,1.3) and (-1.5,1.3) .. (-3,3);%
 \fill[B5Gray] (-1.3,4) circle (.3);\fill[B5Gray] (1.5,4) circle (.3);%
 \draw[b5 link] (0,3.4)--(0,2.7);\draw[b5 link] (-2,-2)--(-2,-5);\draw[b5 link] (2,-2)--(2,-5);%
 \draw[draw=B5Gray,line width=.8pt] (-2.5,-3).. controls (-8,-5) and (-8,-.5) .. (-5,-1);%
\end{scope}}%
\providecommand{\BasicsFiveRobot}[3]{\begin{scope}[shift={(#1,#2)},scale=#3]%
 \fill[B5Gray,rounded corners=.4mm] (-2.7,-1.6) rectangle (-1.7,1.6);%
 \fill[B5Gray,rounded corners=.4mm] (1.7,-1.6) rectangle (2.7,1.6);%
 \draw[draw=B5Gray,fill=white,line width=1pt,rounded corners=.7mm] (-1.8,-2.3) rectangle (1.8,2.3);%
 \draw[draw=B5Gray,fill=B5Yellow,line width=.8pt] (0,.6) circle (.65);%
\end{scope}}%
\BasicsFiveSetup%
\begin{tikzpicture}[b5 picture]%
 \path[use as bounding box] (0,-1) rectangle (111,58);%
 \draw[draw=B5Gray,fill=B5Ice,line width=1pt,rounded corners=1mm] (1,15)--(1,43)--(20,43)--(25,38)--(25,15)--cycle;%
 \draw[b5 link] (20,43)--(20,38)--(25,38);%
 \node[align=center] at (13,29) {我的星海\\订单明天\\能送到\\北京吗？};%
 \coordinate (fork) at (32,29);\draw[b5 link] (25,29)--(fork);\fill[B5Gray] (fork) circle (.65);%
 \foreach \yy/\label in {48/意图识别,29/信息抽取,10/答案生成}{%
  \node[b5 node,minimum width=35mm,minimum height=13mm] (m\yy) at (58,\yy) {};%
  \draw[b5 link] (56.5,\yy-6.5)--(56.5,\yy+6.5);%
  \BasicsFiveNet{43}{\yy}{1.1}%
  \node[font=\figurefont\mdseries\fontsize{7.8}{9.6}\selectfont] at (65.5,\yy) {\label};%
  \draw[b5 flow] (m\yy.east)--(79,\yy);}%
 \draw[b5 flow] (fork) .. controls (38,29) and (35,48) .. (40.5,48);%
 \draw[b5 flow] (fork)--(40.5,29);%
 \draw[b5 flow] (fork) .. controls (38,29) and (35,10) .. (40.5,10);%
 \node[b5 node,minimum width=30mm,minimum height=10mm,fill=B5Ice] at (94,48) {物流查询};%
 \draw[draw=B5Gray,fill=B5Rose,line width=1pt,rounded corners=1mm] (79,22)--(79,36)--(104,36)--(109,31)--(109,22)--cycle;%
 \draw[b5 link] (104,36)--(104,31)--(109,31);%
 \node[align=left] at (94,29) {星海：商家\\北京：地点};%
 \draw[draw=B5Gray,fill=B5Cream,line width=1pt,rounded corners=1.5mm] (79,4)--(79,16)--(109,16)--(109,4)--(85,4)--(81,1)--(82,4)--cycle;%
 \node[align=center] at (94,10) {预计明日\\送达北京。};%
\end{tikzpicture}%
\endgroup%
